\section{Conclusion and Future Work}
\label{sec:conclusion}

\IEEEPARstart{I}{n} this paper, we have evaluated the potential of using GPUs in the OPM project, which contains open-source implementations of several reservoir simulation CFD applications. To perform this study, we have developed several manual ILU0 preconditioned BiCGStab iterative solvers using OpenCL and CUDA and experimented with both AMD and NVIDIA GPUs. Furthermore, we have developed a bridge to seamlessly connect the GPU solvers and third-party scientific libraries into OPM. Finally, we provided extensive evaluation results for all different solver configurations used running on different GPU hardware offered by different hardware vendors. To perform our bench-marking, we use small, medium, and large use cases, starting with the public test case NORNE, which includes approximately 50k active cells, and ending with a large model that includes approximately 1 million active cells. We find that a GPU can accelerate a single dual-thread MPI process up to 5.6 times and that it can compare with around 8 dual-thread MPI processes. 
Although the results are not impressive with respect to GPU performance in other fields that deal with structured computations with regular memory accesses, where massive speedups are observed, the unstructured and irregular nature of the computations in three-phase black oil reservoir simulations makes the GPU performance less impressive. Nevertheless, in this paper, we show that with well-suited techniques, we can improve the performance as it was shown when we compared it to the out-of-the-box library rocsparse \cite{rocsparse}. 
Finally, the source codes for the GPU kernels, solvers, and the integration with OPM \textit{Flow} are made available in the git repositories of the OPM project \cite{opmrepo}. 

In the future, we plan to further analyze the potential slowdowns in the GPU solvers that prevented us from obtaining more impressive speedup numbers. First and foremost, we need to continue with the profiling of the kernels and understand if the matrix data can be better encoded to maximize the cache hit and increase the bandwidth utilized. Furthermore, we should also investigate the impact of adding the well contribution to the matrix or the impact of the memory transfers when the contributions are calculated on the CPU. This could clarify why the speedup for the largest model is smaller than for the small-scale NORNE test case. Finally, several optimizations could be performed, such as using pinned memory in OPM for the assembly part to avoid the implicit CPU-to-CPU memory transfers of the system matrix before the data is copied to the GPU, and overlapping the copy of the original matrix to the GPU while the jacobi matrix is created on the CPU. 
